\documentclass[10pt]{article}

\usepackage[utf8]{inputenc}

\usepackage[T1]{fontenc}

\usepackage{amsmath}

\usepackage{amsfonts}

\usepackage{amssymb}

\usepackage[version=4]{mhchem}

\usepackage{stmaryrd}

\usepackage{bbold}



\begin{document}

Пусть $\gamma$ - произвольная петля в базе расслоения Милнора версальной деформации критич. точки. Ограничение расслоения на $\gamma$ определяет расслоение над окружностью. Соответствующее преобразование монодромии зависит только от гомотопич. типа $\gamma$. Возникает представление фундаментальной группы базы в группу автоморфизмов гомологий слоя, определяюшее М. г.



МОНОДРОМИИ ТЕОРЕМА - достаточныЙ признак однозначности ветви аналитической функции. Пусть $D$ - односвязная область комплексного числового пространства $\mathbb{C}^{n}, n \geqslant 1$. Тогда, если нек-рый элемент аналитич. функции $\Sigma\left(z^{0} ; r\right)$ с центром $z^{0} \in D$ аналитически продолжаем вдоль любого пути, расположенного в $D$, то возникающая при этом аналитич. продолжении ветвь аналитич. функции $f(z)$, $z=\left(z_{1}, \ldots, z_{n}\right)$, однозначна в $D$. Иначе говоря, ветвь аналитич. функции $f(z)$, определяемая односвязной областью $D$ и элементом $\Sigma\left(z^{0} ; r\right)$ с центром $z^{0} \in D$, обязательно однозначна. Другая равносильная формулировка: если элемент $\Sigma\left(z^{0} ; r\right)$ аналитически продолжается вдоль всех путей, принадлежащих произвольной области $D \subset \mathbb{C}^{n}$, то результат этого продолжения в любую точку $z^{*} \in D$ (то есть элемент аналитич. функции $\Sigma\left(z^{*} ; r^{*}\right)$ с центром $\left.z^{*}\right)$ один и тот же для всех гомотопных путей в $D$, соединяющих точки $z^{0}$ и $z^{*}$.



М. т. справедлива и для аналитич. функций $f(z)$, определенных в областях $D$ на римановых поверхностях или на римановых областях.



Лит.: [1] М а р к у ш е в и ч А. И., Теория аналитических функций, 2 изд., т. 2, М., 1968; [2] Владими ров В. С., Методы теории функций многих комплексных переменных, М., 1964.



Е.Д. Соломенцев.



МОНОДРОМИИ ТЕОРИИ ОБРАТНЫЕ ЗАДАЧИ - ЗаДачи восстановления системы обыкновенных дифференциальных уравнений с рациональными коэффициентами:



\[

\left.\begin{array}{c}

\frac{d \Psi}{d \lambda}=A(\lambda) \Psi, \\

\lambda \in \mathbb{C}, \quad A(\lambda)=\sum_{v=1}^{n} \sum_{k=0}^{r_{v}} \frac{A_{v, k}}{\left(\lambda-a_{v}\right)^{k+1}}+\sum_{k=1}^{r_{\infty}} A_{\infty, k} \lambda^{k-1},

\end{array}\right\}

\]



по заданной монодромии группе (п роблема РиманаГ и л ь б е р т а) (см. [1], [2]); здесь $A_{v, k}-(m \times m)$-постоянные по $\lambda$ матрицы $(m>1), a_{1}, \ldots, a_{n}, \infty$ суть $n+1$ различных фиксированных точек на римановой сфере $C P^{1}$.



Группа монодромии задает преобразование фундаментального матричного решения $\Psi$ при обходе вокруг каждой особой точки $a_{v}$ :



\[

\left(\lambda-a_{v}\right) \mapsto\left(\lambda-a_{v}\right) \exp (2 \pi i) \Rightarrow \Psi \mapsto \Psi M_{v} .

\]



Матрицы $M_{\mathrm{v}}$ называются матр ицам и монодроми и. Последовательному обходу контуров, не проходящих через особые точки $a_{v}$, соответствует произведение матриц монодромии, зависящее только от гомотопич. класса пути обхода. Группа монодромии, составленная из таких произведений, является подгруппой $G L(m, C)$, действующей справа на пространстве фундаментальных решений системы (1). Для систем (1) фуксова типа $\left(A_{\infty, k}=0, r_{v}=0\right)$ М.т.о.з. составляет содержание 21-й проблемы Гильберта (см. [3]). Для таких систем эта задача решена для произвольной области $\Omega \subset C$, конформно эквивалентной кругу (см. [4]). В случае $C P^{1}$ решение задачи в классе фуксовых систем всегда существует при $n, m \leqslant 3$. При $n>3$ и любом $m \geqslant 3$ найдется группа монодромии, для к-рой не существует реализующей ее фуксовой системы (см. [5]). Вычисление коэффициентов $A_{v, k}$ в виде бесконечных рядов приведено в [6].



Группа монодромии систем с иррегулярной особой точкой $\left(r_{\infty} \geqslant 1\right)$ содержит дополнительные образующие, называемые матрицам и Стокса и определяемые следующим образом. Пусть $A_{\infty}, r_{\infty}$ имеет не совпадающие между собой собственные числа $\mu_{\alpha}$. Покроем окрестность беско- нечности конечным числом секторов $\Omega_{k}, k=1,2, \ldots, 2 r_{\infty}+1$, так что $\Omega_{k} \cap \Omega_{k+1} \neq \varnothing$ и каждый сектор содержит только один луч раздела, заданный условием $\operatorname{Re}\left[\left(\mu_{\alpha}-\mu_{\beta}\right) \lambda^{r_{\infty}}\right]=0$, $\alpha, \beta=1,2, \ldots, m$. В каждом секторе найдется голоморфное невырожденное решение $\Psi_{k}$ системы (1) с асимптотикой



\[

\Psi_{k}(\lambda)=\left(I+O\left(\lambda^{-1}\right)\right) \exp \left(\sum_{j=1}^{r_{\infty}} D_{j} \lambda^{j}+D_{0} \log \lambda\right), \quad \lambda \rightarrow \infty,

\]



где $D_{j}=\operatorname{diag}\left(d_{1 j}, \ldots, d_{m j}\right), j=0, \ldots, r_{\infty}$. Матрицы Стокса $S_{k}$ связывают решения $\Psi_{k}$ в пересечении соседних секторов:



\[

\Psi_{k+1}(\lambda)=\Psi_{k}(\lambda) S_{k}, k=1,2, \ldots, 2 r_{\infty} .

\]



Для систем с единственной особенностью иррегулярного типа на бесконечности М.т.о. з. разрешима в области $C P^{1} \backslash\{0\}$. Именно, при $|\lambda|>\lambda_{0}$ существует система (1) с матрицей $A$, имеющей вид:



\[

A(\lambda)=\sum_{k=1}^{r_{\infty}} A_{\infty, k^{\lambda}} \lambda^{k-1}+\frac{A_{0}}{\lambda},

\]



такая, что соответствующие матрицы Стокса и матрицы $D_{j}$ совпадают с заданными (см. [2]). Подобно случаю фуксовых систем, решение М.т.о. з. на всем $C P^{1}$ сводится к выделению условий на группу монодромии, при к-рых $A_{0}=0$ (см. [7]). Эти условия нетривиальны даже в случае $m=2$, $r_{\infty}=3$ и связаны с распределением полюсов решений нелинейных уравнений типа Пенлеве, описывающих изомонодромные деформации системы (1) (см. Изомонодромной деформации метод).



Jum.: [1] Fuchs L., Gesammelte mathematische Werke, Bd 1-3, 1904; [2] Birkh off G. D., «Proc. Amer. Acad. Arts and Sci.», 1913, v. 49, n. 9; [3] Проблемы Гильберта, М., 1969; [4] А р н ольд В. И., Ил ья ш е н ко Ю. С., в кн.: Итоги науки и техники. Современные проблемы математики. Фундаментальные направления, М., 1986; [5] Б о л и бр у х А. А., Проблема Римана - Гильберта, «Успехи математич. наук», 1990, т. 45, в. 2, с. 3-47; [6] Е руг ин Н. П., Проблема Римана, Минск, 1976; [7] И тс А. Р., Н о в о кше но в В. Ю., «Функциональный анализ и его прилож.», 1988, т. 22, в. 3, с. 25-36.



В. Ю. Новокшенов.



МОНОДРОМИЯ - автоморфизм гомологий слоя локально тривиального расслоения $\pi: X \rightarrow S^{1}$ такого, что $\left.\pi\right|_{\partial X}-$ тривиальное расслоение. Пусть $V_{t}-$ слой над точкой $t \in S^{1}$. Обход окружности $S^{1}$ с началом и концом в $t_{0}$ поднимается до непрерывного семейства гомеоморфизмов $h_{t}: V_{t_{0}} \rightarrow V_{t}$ такого, что id $=h_{0}: V_{t_{0}} \rightarrow V_{t_{0}} ; h_{2 \pi} \mid \partial V_{t_{0}}=$ id. Линейный оператор $h_{*}: H_{*}\left(V_{t_{0}}\right) \rightarrow H_{*}\left(V_{t_{0}}\right)$, индуцированный $h_{2 \pi}$ и называется монод ромией.



С. В. Чмутов.



MOHOполь - гипотетический магнитный заряд. Экспериментально пока не обнаружен. Одни исследователи вообще отвергают сушествование элементарных магнитных зарядов на том основании, что электродинамика построена в предположении отсутствия источников электромагнитного поля иных, кроме электрич. заряда, и взаимодействие с электромагнитным полем осуществляется только через них (это проявляется в отсутствии дуальной симметрии Максвелла уравнений, то есть симметрии между электрической и магнитной компонентами тензора поля). Другие исследователи, наоборот, полагают, что законы классич. электродинамики не содержат никаких запретов, для к-рых М. не может быть, и допускают существование частиц с одним магнитным полюсом, давая для них определенные уравнения поля и уравнения движения.



Так, напрашивается следуюшее симметричное обобщя ние уравнений Максвелла:



\[

\begin{gathered}

\operatorname{rot} \boldsymbol{B}=4 \pi j+\frac{\partial \boldsymbol{E}}{\partial t}, \quad \operatorname{rot} \boldsymbol{E}=-4 \pi \tilde{j}-\frac{\partial \boldsymbol{B}}{\partial t}, \\

\operatorname{div} \boldsymbol{E}=4 \pi \rho, \quad \operatorname{div} \boldsymbol{B}=4 \pi \tilde{\rho},

\end{gathered}

\]





\end{document}